\section{刑法论 - 刑法的效力}
\begin{enumerate}
  \item 做题第 1 步先判断境内/境外 \yjt \bs{属地/属人、保护、普遍}。
  \item 空间效力。
        \begin{enumerate}
          \item 境内。
                \begin{enumerate}
                  \item \zx{旗国主义}（空海）。例外。\al{国际长途列车、汽车 \yjt 刑诉}。
                  \item 犯罪行为地、结果地。均扩大解释。\zx{预备、教唆、帮助、实行行为；危害、实害、共同犯罪结果}。
                  \item 教唆行为和实行行为的空间割裂。一有一无，\zd{看实行行为是否犯罪}。没考过。
                  \item \zx{单纯过境地}不属于一种紧迫的危险的结果地。
                \end{enumerate}
          \item 境外。
                \begin{enumerate}
                  \item \bs{属人。中国人在国外犯罪}，公职军人一律追究。
                        \begin{enumerate}
                          \item 优先属地。方便侦查。
                          \item 境外贪官是属地。贪污发生在国内，只不过是逃到了国外。
                        \end{enumerate}
                  \item 保护。对中国公民实施犯罪 + \zx{3 年以上+ 双重犯罪}原则。\al{章莹颖}。
                  \item 普遍。针对国际犯罪。\zd{海盗、种族灭绝、恐怖犯罪}。考的少。
                \end{enumerate}
        \end{enumerate}
  \item 时间效力。
        \begin{enumerate}
          \item \zx{再审}案件只能\zx{从旧}，适用行为时。\al{再审流氓罪}。
          \item 司法解释的时间效力。
                \begin{enumerate}
                  \item \bs{行无审有}。行为时无，审判时有 \yjt 按司法解释。
                  \item \bs{行有审新，从旧兼从轻（仅针对未决案）}。行为时有，审判时有新的司法解释 \yjt 从旧兼从轻。
                \end{enumerate}
        \end{enumerate}
\end{enumerate}
